% Generated draft: requires mathematical and source review.
\documentclass[11pt]{article}
\usepackage[margin=25mm]{geometry}
\usepackage{fontspec}
\setmainfont{Latin Modern Roman}
\usepackage{amsmath,amssymb,mathtools,unicode-math}
\setmathfont{Latin Modern Math}
\usepackage{enumitem,needspace,xurl,hyperref}
\setlength{\parindent}{0pt}
\setlength{\parskip}{5pt}
\setlength{\emergencystretch}{3em}
\newcommand{\N}{\mathbb{N}}
\newcommand{\Z}{\mathbb{Z}}
\newcommand{\Q}{\mathbb{Q}}
\newcommand{\R}{\mathbb{R}}
\newcommand{\C}{\mathbb{C}}
\newcommand{\F}{\mathbb{F}}
\newcommand{\A}{\mathbb{A}}
\newcommand{\PP}{\mathbb{P}}
\newcommand{\E}{\mathbb{E}}
\newcommand{\eps}{\varepsilon}
\newcommand{\dd}{\,\mathrm{d}}
\begin{document}
\section*{UnsolvedMath ID: 30004464}
\textbf{Source problem code:} catalogue entry
\subsection{Definitions and mathematical statement}
The supplied formulation leaves the group category, the meaning of smoothness, and the topology on the subalgebra unspecified. The following is a conditional formulation under the usual discrete-group interpretation; these choices require confirmation from the cited report.

Let \(\Gamma\) be a discrete group. Its complex group algebra \(\C\Gamma\) consists of finite sums \(\sum_{g\in\Gamma}c_g g\), with multiplication induced by the group law. Let \(\lambda(g)\xi(h)=\xi(g^{-1}h)\) for \(\xi\in\ell^2(\Gamma)\), and define the reduced group \(C^*\)-algebra by
\[
C_r^*(\Gamma)=\overline{\lambda(\C\Gamma)}^{\|\cdot\|}\subseteq\mathcal B(\ell^2(\Gamma)).
\]
Identify \(\C\Gamma\) with its image under \(\lambda\). To formulate an extension by restriction, assume explicitly that
\[
\C\Gamma\subseteq A\subseteq C_r^*(\Gamma),
\]
where \(A\) is a complex subalgebra dense in the \(C^*\)-norm and is equipped with a specified locally convex algebra topology \(\mathcal T\). Neither this topology nor the additional properties intended by ``suitable smooth'' are supplied; norm density alone does not specify them.

For an integer \(n\geq0\), an algebraic cyclic \(n\)-cocycle on \(\C\Gamma\) means an \((n+1)\)-linear map \(\varphi:(\C\Gamma)^{n+1}\to\C\) satisfying
\[
\varphi(a_0,\ldots,a_n)=(-1)^n\varphi(a_n,a_0,\ldots,a_{n-1})
\]
and
\[
\begin{aligned}
(b\varphi)(a_0,\ldots,a_{n+1})
&=\sum_{j=0}^{n}(-1)^j\varphi(a_0,\ldots,a_ja_{j+1},\ldots,a_{n+1})\\
&\quad+(-1)^{n+1}\varphi(a_{n+1}a_0,a_1,\ldots,a_n)=0
\end{aligned}
\]
for every \(a_0,\ldots,a_{n+1}\in\C\Gamma\).

\textbf{Extension question.} Determine conditions on \(\Gamma\), \((A,\mathcal T)\), \(n\), and \(\varphi\) under which there exists a jointly continuous \((n+1)\)-linear map \(\widetilde\varphi:A^{n+1}\to\C\) satisfying the same cyclicity and cocycle identities for all arguments in \(A\), and such that
\[
\widetilde\varphi\big|_{(\C\Gamma)^{n+1}}=\varphi.
\]
Here joint continuity is the explicit convention used in this conditional formulation. The source asks for extension of a cocycle itself, rather than merely extension of its cyclic cohomology class. It gives a request for conditions, not an assertion that every cocycle extends to every smooth dense subalgebra.

In degree zero the cocycle identity is the trace identity \(\varphi(ab)=\varphi(ba)\). The supplied dated assessment mentions partial answers for conjugacy-class traces under a sufficiently large spectral-gap hypothesis, but supplies neither the precise traces covered nor the operator, gap bound, or other hypotheses. Thus no precise partial extension theorem can be stated from this record alone.
\subsection{Short English statement}
Determine when a cyclic cocycle on a group algebra extends continuously, with its values preserved, to a smooth dense subalgebra of the reduced group \(C^*\)-algebra. A precise instance requires the group category, containment of the group algebra, the topology, and the meaning of smoothness to be specified. The discrete-group formulation above is conditional. The supplied partial-solution assessment requires independent review.
\subsection{Sources}
\begin{itemize}
\item Local UnsolvedMath registry: supplied complete record, permanent ID \(30004464\), title ``Extension of Group Cyclic Cocycles to Smooth Subalgebras,'' Oberwolfach Reports Open Problems collection, exact source problem code \texttt{catalogue entry}. No individual proposer is identified; the supplied proposed year is 2020.
\item Supplied source citation and literature source labeled ``original OWR report'': \emph{Non-Commutative Geometry and Cyclic Homology} (2021), \emph{Oberwolfach Reports}, pp.\ 919--960, DOI \texttt{10.4171/owr/2020/17}, \url{https://doi.org/10.4171/owr/2020/17}. The background attributes the collection to workshop reports and their open-problem discussions. This report reference is supplied provenance and has not been independently consulted here.
\item Supplied literature assessment, checked on \textbf{2026-08-22}: the record labels the problem ``partially solved,'' describes the general-group question as still open, and reports partial answers for degree-zero conjugacy-class traces under a sufficiently large spectral-gap hypothesis. This is supplied dated curation, not independently verified literature status or a verified theorem. The background explicitly requests independent verification of time-sensitive status and cited claims.
\item The supplied statement-status label is \texttt{corrected\_verified}; its accompanying explanation says that mathematical content was preserved while notation, typography, omitted definitions, or duplicate clauses were repaired. This records the registry's assessment and does not constitute independent mathematical verification of this draft.
\end{itemize}
\end{document}
